\begin{definition}[Subgrupo Normal]
  Um subgrupo $N$ de um grupo $G$ é \textbf{normal em $G$} se $gN=Ng$, para todo $g\in G$. Neste caso, denotamos $N\normal G$.  
\end{definition}
\begin{proposition}
  Um subgrupo $N$ de $G$ é normal em $G$ se, e somente se, $gag^{-1}\in N$, para todo $g\in G$ e todo $a\in N$.   
\end{proposition}
\begin{example}
  $\{e\}$ e $G$ são normais em $G$. 
\end{example}
\begin{example}
  Se $G$ é umm grupo abeliano, então todos os seus subgrupos são normais. No entanto a recíproca não e verdadeira, isto é, se todos os subgrupos de um grupo $G$ são normais, então $G$ não é necessarimente abeliano, pode tomar como exemplo $S_3$.
\end{example}
\begin{example}
  $Z(G)\normal G$. Mais geralmente, se $H\leq Z(G)$ então $H\normal G$. 
\end{example}
\begin{proposition}
  Se $N$ é subgrupo normal de $G$ então o conjunto $G/N=\{aN:a\in G\}$ é um grupo com a operação $Na\cdot Nb=Nab$, $a,b\in G$. O grupo $G/N$ é dito o \text{grupo quociente} de $G$ por $N$.

  Na verdade, $G/N$ é um grupo se, somente se, $N\normal G$. 
\end{proposition}
